% ======================================================================
% 第4章 関連研究（経路制御における既存研究と課題）
% ======================================================================

% 【追加】章の導入（位置付け）
本章では，第2章で定義した多制約最適経路問題 (MCOP) に対する既存の解決手法を概観し，それぞれの特徴と課題について整理する．
まず，厳密解法の計算量的限界について述べ，次に，遺伝的アルゴリズム (GA) や蟻コロニー最適化 (ACO) といった他のメタヒューリスティクスの適用事例とその課題を挙げる．
最後に，本研究の基礎となるPSOを用いた経路制御の先行研究について解説し，既存手法では解決できていない課題（早熟収束や計算時間）を明確にすることで，第5章で提案する手法の必要性と本研究の立ち位置を明らかにする．

\section{厳密解法とその限界}
多制約経路問題 (MCOP) に対する厳密解法としては，SAMCRA (Self-Adaptive Multiple Constraints Routing Algorithm)\cite{van2004} やラベル修正法 (Label Correcting Method) \cite{guerriero2001}が代表的である．
Van Mieghemらは，非線形な長さ関数と$k$-shortest pathアルゴリズムを組み合わせることで，探索空間を効率的に枝刈りするSAMCRAを提案した\cite{van2004}．
なお，本研究が対象とするボトルネックリンク問題 (MBL) 単体に対しては，Weiらによる改良型ダイクストラ法\cite{Wei2019}や，Malpaniらによる帯域幅最大化パス構築法\cite{Malpani2002}などが提案されている．
しかし，これらに遅延制約を加えたMCOPとなると，前述の通りNP困難性が壁となり，ネットワーク規模の増大に伴い計算時間が指数関数的に増加するという致命的な課題がある．
特に，ノード数が数千を超える大規模ネットワークや，制約条件が緩く候補パスが大量に存在する場合，実用的な時間内すなわち数ミリ秒オーダーで解を出力することは困難である．
したがって，リアルタイム性が求められる次世代ネットワークの制御においては，厳密性よりも応答速度を重視した近似解法が不可欠となる．

\section{他のメタヒューリスティクスによるアプローチ}
厳密解法の計算量問題を解決するため，遺伝的アルゴリズム (GA) や蟻コロニー最適化 (ACO) といったメタヒューリスティクスの適用が研究されてきた．

GAを用いたアプローチとして，Moshrefらは無線センサネットワーク (WSN) におけるQoS最適化を目的としたENSGRAを提案している\cite{Moshref2021}．彼らは多親交叉 (Multi-parent Crossover) を導入することでパレート最適解の探索能力を強化し，多目的PSOと比較して高い最適化性能を達成した．
また，Nairらは帯域幅最大化を目的としたGAにおいて，突然変異操作がリンク切断などの不連続な探索空間からの脱出に有効であることを示している\cite{Nair2014}．
しかし，GAは交叉や突然変異によって生成された経路が物理的に接続されていない（非実行可能解となる）場合が多く，その修復処理やペナルティ計算に多大な計算コストを要するため，リアルタイム制御には不向きである．

一方，ACOを用いたアプローチとして，Supiyandiらはネットワーク最適化における比較実装を行い，高トラフィック環境下においてACOがPSOよりも優れたスループットと低いパケット損失率を維持することを示した\cite{Supiyandi2024}．ACOはフェロモンを用いた正のフィードバックにより，混雑状況の変化に対して自律的に迂回経路を選択できる強みを持つ．
しかし，探索初期はフェロモンが蓄積されていないため収束に時間を要するコールドスタート問題や，多数の探索エージェントを管理するためのオーバーヘッドが大きく，大規模ネットワークでの高速な応答性には課題が残る．

これらの単独手法の課題に対し，PatelらはACOとPSOを組み合わせることで，遅延や帯域などのQoS制約を満たすマルチキャストルーティング手法を提案している\cite{Patel2014}．
このようなハイブリッド手法は高い探索能力を示す一方で，アルゴリズムが複雑化し，調整すべきパラメータが増加するという新たな課題を生む．

\section{PSOを用いた経路制御手法}
PSOは，GAやACOと比較してパラメータが少なく，実装が容易であり収束が速いことから，近年QoSルーティングへの適用が進んでいる．

PSOを経路制御に適用する先駆的な研究として，Mohemmedらは各ノードに優先度を割り当てる間接エンコーディング手法を提案し，最短経路問題においてその有効性を示した\cite{Mohemmed2008}．
また，本庄らは，巡回セールスマン問題 (TSP:Traveling Salesman Problem) および経路制御に対しPSOを適用し，近傍構造すなわちトポロジーの影響を分析した\cite{honjo2013}．
彼らの研究では，全粒子が情報を共有するgBestモデルよりも，近傍のみと情報を共有するlBestモデルの方が，解の多様性が維持され局所解に陥りにくいことが示されている．

本研究の直接的な先行研究として，乗松らはボトルネックリンク帯域最大化 (MBL) 問題に対してPSOを適用する手法を提案している\cite{norimatsu2024}．
彼らの研究では，各ノードに実数値の優先度を付与し，その優先度に基づいて経路を構築するエンコーディング手法を用いることで，ループのない実行可能な経路を効率的に生成できることを示した．また，シミュレーション実験により，提案手法が従来のダイクストラ法と比較して，より高いボトルネック帯域を持つ経路を発見できることを実証している．
しかし，彼らの手法は帯域幅のみを考慮した単一目的最適化にとどまっており，遅延などの加法的な制約を含むMCOPへの拡張は行われていない．また，探索の停滞に対するリスタート機構も実装されておらず，より複雑な制約条件下や大規模ネットワークにおける探索性能については課題が残されていた．

\section{本研究の位置付け}
以上の関連研究の調査に基づき，本研究の独自性と優位性は以下の3点に集約される．

第一に，lBestトポロジーを採用する．
既存の多くのQoSルーティングが採用しているgBestモデルではなく，局所解回避能力の高いlBestモデルを採用することで，MCOPの複雑な制約条件に対応する．

第二に，リスタート戦略を導入する．
乗松らの先行研究や他の既存手法で不足していた停滞からの脱出策として，群れの分散状況に応じた動的リスタート戦略を導入し，ハイブリッド手法のような複雑さを招くことなく，ACOなどよりも高速な環境適応を実現する．

第三に，並列化による高速化を行う．
lBestモデルの欠点である収束速度の遅さを補うため，マルチコアプロセッサによる並列実装を行い，厳密解法では不可能な大規模ネットワークでのリアルタイム制御を実現する．

すなわち，本研究はアルゴリズムの改良と実装技術の両面から，次世代ネットワークに求められる高速かつロバストな経路制御手法を確立するものである．